% ==============================================================================
% APPENDIX.TEX - PHỤ LỤC
% ==============================================================================

\chapter*{PHỤ LỤC}
\addcontentsline{toc}{chapter}{PHỤ LỤC}
\setcounter{section}{0}
\renewcommand{\thesection}{\Alph{section}}
\renewcommand{\theHsection}{appendix.\Alph{section}}

% Macro định nghĩa từng mục phụ lục: không dùng \numberline trong TOC để tránh lỗi đè chữ
\newcommand{\appendixsection}[2]{%
    \refstepcounter{section}%
    \section*{Phụ lục \Alph{section}: #1}%
    \addcontentsline{toc}{section}{Phụ lục \Alph{section}: #1}%
    \label{#2}%
}

\appendixsection{Đặc tả kỹ thuật Dockerfile cho dịch vụ AI Serverless}{app:dockerfile}
\begin{lstlisting}[language=bash, caption={Cấu hình Dockerfile tối ưu kích thước cho Google Cloud Run}]
# Build stage: Base Python 3.11 Slim
FROM python:3.11-slim

# Thiet lap bien moi truong
ENV PYTHONUNBUFFERED=1 \
    PYTHONDONTWRITEBYTECODE=1 \
    PORT=8080

# Cai dat thu vien he thong cho OpenCV Headless va MediaPipe
RUN apt-get update && apt-get install -y --no-install-recommends \
    libglib2.0-0 \
    && rm -rf /var/lib/apt/lists/*

WORKDIR /app

# Cai dat Python dependencies
COPY requirements.txt .
RUN pip install --no-cache-dir -r requirements.txt

# Sao chep ma nguon va trong so mo hinh
COPY src/ /app/src/
COPY train/vsl_cnn_pose_keras.h5 /app/train/vsl_cnn_pose_keras.h5

# Cau hinh tai khoan non-root bao mat
RUN useradd -m -u 1000 appuser && chown -R appuser:appuser /app
USER appuser

# Expose cong va chay Uvicorn server
EXPOSE 8080
CMD ["uvicorn", "src.vsl.service:app", "--host", "0.0.0.0", "--port", "8080"]
\end{lstlisting}

\appendixsection{Cấu hình luồng CI/CD tự động hóa trên GitHub Actions}{app:github_actions}
\begin{lstlisting}[language=bash, caption={Workflow CI tự động chạy kiểm thử trên GitHub Actions (.github/workflows/ci.yml)}]
name: CI Pipeline - Lint and Automated Tests

on:
  push:
    branches: [ develop, main ]
  pull_request:
    branches: [ develop, main ]

jobs:
  test:
    runs-on: ubuntu-latest
    steps:
      - name: Checkout Source Code
        uses: actions/checkout@v4
        with:
          lfs: true

      - name: Setup Python Runtime
        uses: actions/setup-python@v5
        with:
          python-version: '3.12'
          cache: 'pip'

      - name: Install Dependencies
        run: |
          python -m pip install --upgrade pip
          pip install -r requirements.txt
          pip install -r requirements-dev.txt

      - name: Run Test Suite
        run: |
          pytest tests -v
\end{lstlisting}
